\chapter{Quantum Foundations and the Measurement Problem}
\label{ch:quantum-foundations}

\epigraph{\itshape I think I can safely say that nobody understands
quantum mechanics.}{--- Richard Feynman, \textit{The Character of
Physical Law}, 1965}

\noindent
This chapter develops the quantum-mechanical apparatus we will need
in subsequent chapters. The treatment is pedagogical: it assumes no
prior physics. The reader who has worked through an undergraduate
physics course will find the early sections elementary; the reader
who has not will find them sufficient. The mathematical level is
calibrated to the conceptual demands of the subsequent material;
where we need a piece of mathematics, we develop it.

The chapter has a specific purpose within the larger work. We are
not writing a quantum-mechanics textbook. We are extracting the
specific features of the quantum-mechanical formalism that bear on
the questions Part III will address: the relationship between
determination and freedom (\arb{qadar} and \arb{ikhtiyar}), the
status of probability in non-classical systems, the role of the
observer in instantiating outcomes, and the philosophical commitments
that different interpretations of the formalism imply. Each of these
will be developed only to the depth required.

The chapter proceeds in nine sections.
\xref{sec:qm-why} explains why a chapter on quantum mechanics is in
this book at all.
\xref{sec:complex-vectors} develops the mathematical preliminaries
(complex numbers, vectors, Hilbert spaces) for readers who need them.
\xref{sec:wavefunction} introduces the wave function and the
principle of superposition.
\xref{sec:schrodinger} treats the Schr\"odinger equation and unitary
evolution.
\xref{sec:born-rule} treats observables, eigenvalues, and the Born
rule.
\xref{sec:measurement} states the measurement problem.
\xref{sec:interpretations} surveys the major interpretations
(Copenhagen, Many-Worlds, Bohmian mechanics, QBism, objective
collapse, relational quantum mechanics) with their philosophical
commitments.
\xref{sec:bell} treats Bell's theorem and the question of locality.
\xref{sec:qm-implications} draws the implications for the Hidden
Architecture project, particularly for the treatment of \arb{qadar}
in Chapter 16.

\section{Why a chapter on quantum mechanics is in this book}
\label{sec:qm-why}

The reader may reasonably ask why an Islamic economics treatise
contains a quantum-mechanics chapter. The answer was previewed in
\xref{ch:modern-tools}; we expand on it here.

The classical \arb{kalam} debate on \arb{qadar} (treated in
\xref{sec:kalam}) attempted to reconcile two theological commitments:
that all events are determined by divine decree, and that human
choice is real and morally responsible. The Mu`tazila preserved
human agency by restricting divine creative activity to events
other than human acts. The Ash`ariyya preserved divine omnipotence
by adopting the doctrine of \arb{kasb}, in which God creates human
acts but humans ``acquire'' them. The Maturidiyya took a middle
position. Each formulation has strengths; each leaves residual
puzzles.

The puzzles persist because the classical theologians were working
within a Newtonian-Aristotelian framework in which determinism and
the reality of motion seemed mutually exclusive. If everything is
determined, motion is a mere unfolding of what was already settled;
if motion is real, things could have been otherwise; the two seem
incompatible. The classical theologians did remarkable work
articulating positions within this framework, but the framework
itself constrained what could be said.

The framework began to come apart in the early twentieth century.
The discovery of quantum mechanics --- and, more importantly, the
ongoing dispute about what quantum mechanics means --- has shown
that several of the classical assumptions about determinism, motion,
probability, and the role of the observer were not necessary. Other
options exist. Specifically, two options exist that map cleanly onto
the structure the classical \arb{kalam} debate was trying to
articulate.

The first is \emph{Bohmian mechanics}. In Bohmian mechanics, the
wave function (which is real and deterministic) guides particles
that have definite trajectories at all times. The trajectory of any
given particle is fully determined by the initial wave function and
the particle's initial position. Yet the motion is real: the
particle actually moves, the trajectory is the description of its
motion, and the wave function does not replace the particle but
guides it. Determinism and the reality of motion coexist.

This is exactly the structure that the Ash`ari doctrine of
\arb{kasb} was attempting. God creates the wave function (the
trajectory of the universe); humans ``acquire'' their actions in the
sense that the actions are the actual motion through the
configuration space that the wave function determines. Both
descriptions are correct; neither contradicts the other; the
apparent contradiction was a product of Newtonian assumptions that
quantum mechanics has dissolved.

The second is \emph{QBism} (Quantum Bayesianism). In QBism, the wave
function does not describe an objective state of the world; it
describes the beliefs of an agent. Different agents can have
different valid wave functions for the same system, depending on
what each agent knows. Probability, in QBism, is irreducibly
perspectival.

This too maps onto a classical \arb{kalam} concern. Divine knowledge
is perfect; human knowledge is finite. From the divine perspective,
all outcomes are known; from the human perspective, outcomes are
uncertain and probabilistic. The two perspectives are both correct
because probability is perspectival; there is no contradiction.

These are not idle philosophical curiosities. They are working
interpretations defended by serious physicists, and they provide the
formal vocabulary in which the Islamic position on \arb{qadar} can
be articulated without contradiction. To show this, we must develop
the quantum mechanics carefully enough that the structure becomes
visible. This chapter does that work.

\section{Mathematical preliminaries}
\label{sec:complex-vectors}

Before we can develop quantum mechanics we need three mathematical
preliminaries: complex numbers, vector spaces, and the inner
product. Readers familiar with these may skim or skip the section.

\subsection{Complex numbers}

A \emph{complex number} is a number of the form $z = a + bi$ where
$a$ and $b$ are real numbers and $i$ is a symbol satisfying $i^2 =
-1$. The number $a$ is called the \emph{real part} of $z$, written
$\text{Re}(z)$; the number $b$ is the \emph{imaginary part},
written $\text{Im}(z)$.

Complex numbers obey the standard arithmetic of real numbers, with
the additional rule that $i^2 = -1$:

\begin{align}
(a + bi) + (c + di) &= (a+c) + (b+d)i \\
(a + bi) \cdot (c + di) &= (ac - bd) + (ad + bc)i.
\end{align}

The \emph{complex conjugate} of $z = a + bi$ is $\bar{z} = a - bi$.
The \emph{modulus} (or absolute value) of $z$ is
$|z| = \sqrt{a^2 + b^2} = \sqrt{z \bar{z}}$. The square of the
modulus, $|z|^2 = z\bar{z} = a^2 + b^2$, is real and non-negative;
it will play a major role in what follows.

A complex number can be visualized as a point in a two-dimensional
plane (the \emph{complex plane}), with the real part along the
horizontal axis and the imaginary part along the vertical. The
modulus $|z|$ is the distance from the origin to the point.

Equivalently, every complex number can be written in \emph{polar
form} as $z = r e^{i\theta}$, where $r = |z|$ and $\theta$ is the
angle measured counterclockwise from the positive real axis. Euler's
formula gives the connection: $e^{i\theta} = \cos\theta + i\sin\theta$.

\begin{example}[Multiplication of complex numbers in polar form]
If $z_1 = r_1 e^{i\theta_1}$ and $z_2 = r_2 e^{i\theta_2}$, then
$z_1 z_2 = r_1 r_2 e^{i(\theta_1 + \theta_2)}$. Multiplication
multiplies the moduli and adds the phases. This is structurally
significant: it is why quantum interference, which involves
multiplication of complex amplitudes, depends on the phases as well
as the magnitudes.
\end{example}

\subsection{Vector spaces}

A \emph{vector space over a field} (we will use either the real
numbers $\mathbb{R}$ or the complex numbers $\mathbb{C}$) is a set
$V$ equipped with two operations:

\begin{itemize}[leftmargin=2em]
    \item \emph{Vector addition}: for any $u, v \in V$, the sum
    $u + v \in V$.
    \item \emph{Scalar multiplication}: for any $c$ in the field
    (so $c \in \mathbb{R}$ or $c \in \mathbb{C}$) and any $v \in V$,
    the product $c v \in V$.
\end{itemize}

These operations satisfy the standard rules: addition is
commutative and associative, scalar multiplication distributes over
addition, and so on. The reader has met vectors before in the form
of arrows in space, where addition is geometric (tip-to-tail) and
scalar multiplication is rescaling.

For quantum mechanics, the relevant vectors live in a complex vector
space, meaning the scalars used are complex numbers rather than
real. This generalizes the geometric intuition: vectors can be added
and multiplied by complex scalars to form new vectors in the same
space.

A vector space has a \emph{basis}: a minimal set of vectors
$\{e_1, e_2, \dots, e_n\}$ such that every vector $v$ in the space
can be written uniquely as a linear combination
\[
    v = c_1 e_1 + c_2 e_2 + \dots + c_n e_n
\]
for some scalars $c_1, c_2, \dots, c_n$. The number $n$ is the
\emph{dimension} of the space. For finite-dimensional spaces, $n$
is finite; for infinite-dimensional spaces (which arise in quantum
mechanics), the basis is infinite and the linear combination becomes
an infinite series.

\subsection{Inner products and Hilbert space}

An \emph{inner product} on a complex vector space $V$ is a function
that takes two vectors $u, v \in V$ and returns a complex number,
written $\langle u, v \rangle$ or $\langle u | v \rangle$. The inner
product satisfies:

\begin{itemize}[leftmargin=2em]
    \item \emph{Linearity in the second argument}:
    $\langle u, c_1 v_1 + c_2 v_2 \rangle =
     c_1 \langle u, v_1 \rangle + c_2 \langle u, v_2 \rangle$.
    \item \emph{Conjugate symmetry}:
    $\langle u, v \rangle = \overline{\langle v, u \rangle}$.
    \item \emph{Positive definiteness}: $\langle v, v \rangle \geq 0$
    for all $v$, with equality only if $v = 0$.
\end{itemize}

The inner product permits the definition of \emph{length} (or
\emph{norm}) of a vector: $\|v\| = \sqrt{\langle v, v \rangle}$.
And of \emph{orthogonality}: vectors $u$ and $v$ are orthogonal if
$\langle u, v \rangle = 0$.

A \emph{Hilbert space} is a complex inner-product space that is
``complete'' in a technical sense (every Cauchy sequence converges
to a limit in the space). The technical detail does not concern us;
what matters is that Hilbert spaces are the natural setting for
quantum mechanics.

In quantum mechanics, the inner product is conventionally written
using \emph{Dirac notation}: $\langle \phi | \psi \rangle$ for the
inner product of $\phi$ and $\psi$. The vector $|\psi\rangle$ is
called a \emph{ket}; the vector $\langle\phi|$ is called a \emph{bra};
together they form a ``bra-ket,'' which is the inner product.

\begin{notationbox}[Bra-ket conventions]
\begin{itemize}[leftmargin=2em]
    \item $|\psi\rangle$ is a vector (a ``ket'') in the Hilbert space.
    \item $\langle\phi|$ is the corresponding linear functional (a
    ``bra'').
    \item $\langle\phi|\psi\rangle$ is a complex number, the inner
    product of $\phi$ and $\psi$.
    \item The norm of $|\psi\rangle$ is $\|\psi\| =
    \sqrt{\langle\psi|\psi\rangle}$, which is real and non-negative.
    \item A vector with norm 1 is called \emph{normalized}.
\end{itemize}
\end{notationbox}

\section{The wave function and superposition}
\label{sec:wavefunction}

\subsection{States as vectors}

The first postulate of quantum mechanics:

\begin{principlebox}[Postulate 1: States are vectors in Hilbert space]
The state of a quantum system is described by a normalized vector
$|\psi\rangle$ in a complex Hilbert space $\mathcal{H}$. Two vectors
that differ by a global phase --- $|\psi\rangle$ and $e^{i\theta}
|\psi\rangle$ for any real $\theta$ --- represent the same physical
state.
\end{principlebox}

This is a strong claim and deserves emphasis. The state of a
quantum system --- everything that can be known about it for the
purposes of physics --- is captured by a vector in a complex Hilbert
space. The vector encodes, in a precise sense we will spell out, the
probabilities of every possible outcome of every possible measurement
on the system.

What a vector represents depends on the system. For a single quantum
particle in three-dimensional space, the wave function is a complex-valued
function $\psi(x, y, z)$ on three-dimensional space; the Hilbert space
is the space of all such (square-integrable) functions. For a quantum
spin system, the wave function is a finite-dimensional vector; the
Hilbert space is $\mathbb{C}^2$ for a single spin-1/2 particle, or
$\mathbb{C}^4$ for a pair of them, and so on.

\subsection{Superposition}

The fundamental novelty of quantum mechanics is the
\emph{superposition principle}: if $|\psi_1\rangle$ and $|\psi_2\rangle$
are valid quantum states, then any normalized linear combination
\[
    |\psi\rangle = \alpha |\psi_1\rangle + \beta |\psi_2\rangle, \qquad
    |\alpha|^2 + |\beta|^2 = 1
\]
is also a valid quantum state. The system in state $|\psi\rangle$
is said to be in a \emph{superposition} of $|\psi_1\rangle$ and
$|\psi_2\rangle$ with amplitudes $\alpha$ and $\beta$.

The superposition principle has no analogue in classical physics. A
classical particle can be at position $x_1$ or at position $x_2$,
but not in a superposition of both. A quantum particle can be in a
superposition of being at $x_1$ and being at $x_2$, with the
superposition being a perfectly definite (and physically distinct)
state.

\begin{example}[Schr\"odinger's cat]
The famous thought experiment: a cat in a sealed box is connected to
a quantum coin-flip device. If the device produces outcome $A$, the
cat lives; outcome $B$, the cat dies. Before the box is opened, the
combined system is in state
\[
    |\psi\rangle = \frac{1}{\sqrt{2}} \left( |\text{alive}\rangle +
    |\text{dead}\rangle \right).
\]
The cat, on the standard reading, is in a superposition of being
alive and being dead. Schr\"odinger proposed the example as a
reductio of the Copenhagen interpretation: surely a cat cannot be
in such a superposition? But the mathematics says exactly that, and
the question of how to resolve the apparent paradox is one of the
central problems of quantum interpretation.
\end{example}

\subsection{The complex amplitudes have meaning}

A crucial fact: the complex numbers $\alpha$ and $\beta$ in a
superposition are not just bookkeeping. Their relative phases ---
the difference in their angular position in the complex plane --- can
have observable consequences through interference effects. Two
superpositions with the same magnitudes $|\alpha|$ and $|\beta|$ but
different relative phases describe physically distinct states.

This is the source of quantum interference, the most distinctive
empirical signature of quantum mechanics. The double-slit experiment
--- in which a single particle, sent through a barrier with two
slits, produces an interference pattern characteristic of waves ---
is the classical demonstration. The pattern arises because the
amplitudes for the particle going through one slit and through the
other add coherently, with their relative phases determining where
constructive and destructive interference occur.

\section{The Schr\"odinger equation and unitary evolution}
\label{sec:schrodinger}

\subsection{Time evolution is deterministic}

The second postulate of quantum mechanics:

\begin{principlebox}[Postulate 2: Time evolution is deterministic]
A closed quantum system evolves according to the time-dependent
Schr\"odinger equation:
\[
    i\hbar \frac{d}{dt} |\psi(t)\rangle = \hat{H} |\psi(t)\rangle,
\]
where $\hbar$ is the reduced Planck constant and $\hat{H}$ is the
Hamiltonian operator (the energy operator) of the system.
\end{principlebox}

The Schr\"odinger equation is a deterministic differential equation.
Given a starting state $|\psi(0)\rangle$ and the Hamiltonian
$\hat{H}$, it determines $|\psi(t)\rangle$ for all subsequent times.
There is no randomness in this evolution; it is fully predictable.

Equivalently, we can write the evolution as
\[
    |\psi(t)\rangle = U(t) |\psi(0)\rangle, \qquad U(t) = e^{-i\hat{H}t/\hbar},
\]
where $U(t)$ is a unitary operator (an operator that preserves inner
products and norms). The state at time $t$ is obtained from the
state at time 0 by applying the unitary $U(t)$. Unitary evolution
preserves superposition: if the initial state was a superposition,
so is the evolved state, with the amplitudes evolving in a definite
way.

\subsection{Determinism and unpredictability}

Note carefully what the determinism does and does not say. The wave
function evolves deterministically. But the wave function does not
correspond directly to observable quantities; what is observable is
the result of measurements, which (according to the standard
treatment) introduces probability. The deterministic evolution of
$|\psi\rangle$ is consistent with probabilistic measurement
outcomes.

This is the central tension of standard quantum mechanics. The
mathematics has both a deterministic part (Schr\"odinger evolution)
and a probabilistic part (measurement). These two parts seem to be
of fundamentally different kinds. Most of the interpretive disputes
concern how to reconcile them.

\section{Observables, eigenvalues, and the Born rule}
\label{sec:born-rule}

\subsection{Observables as operators}

The third postulate:

\begin{principlebox}[Postulate 3: Observables are Hermitian operators]
Every physical observable (position, momentum, energy, spin, etc.)
is represented by a Hermitian operator on the Hilbert space. A
Hermitian operator $\hat{A}$ satisfies $\langle\phi | \hat{A} \psi
\rangle = \langle \hat{A}\phi | \psi\rangle$ for all $|\phi\rangle$
and $|\psi\rangle$.
\end{principlebox}

Hermitian operators have two key properties that make them suitable
for representing observables:

\begin{enumerate}[leftmargin=2em]
    \item Their eigenvalues are real numbers.
    \item Their eigenvectors form a complete orthonormal basis for
    the Hilbert space.
\end{enumerate}

Recall that for an operator $\hat{A}$ and a vector $|a\rangle$, if
$\hat{A} |a\rangle = a |a\rangle$ for some scalar $a$, we say
$|a\rangle$ is an \emph{eigenvector} (or \emph{eigenstate}) of
$\hat{A}$ with \emph{eigenvalue} $a$. The eigenvectors of a
Hermitian operator span the Hilbert space, meaning every state can
be written as a superposition of eigenstates.

\subsection{The Born rule}

The fourth postulate is where probability enters quantum mechanics.

\begin{principlebox}[Postulate 4: The Born rule]
When a measurement of observable $\hat{A}$ is performed on a system
in state $|\psi\rangle$, the only possible outcomes are the
eigenvalues of $\hat{A}$. The probability of obtaining eigenvalue
$a$ is
\[
    P(a) = |\langle a | \psi\rangle|^2,
\]
where $|a\rangle$ is the (normalized) eigenvector of $\hat{A}$
corresponding to eigenvalue $a$.

After the measurement, the state of the system collapses to
$|a\rangle$.
\end{principlebox}

The Born rule has three components, each of which is a strong
claim.

First, the only possible outcomes are eigenvalues. Whatever value
the system is in superposition over, the measurement always returns
one of the eigenvalues, never a value in between or outside.

Second, the probability of each outcome is given by the squared
modulus of the corresponding amplitude. This is where the complex
numbers connect to observable probabilities; the squared modulus is
real and non-negative, as a probability should be.

Third, after the measurement, the state collapses to the eigenstate
corresponding to the observed eigenvalue. The system is no longer
in a superposition; it is in a definite eigenstate of the measured
observable.

\subsection{An example: spin}

Consider a spin-1/2 particle, the simplest non-trivial quantum
system. The Hilbert space is $\mathbb{C}^2$, with basis vectors
$|0\rangle$ and $|1\rangle$ (sometimes written $|\uparrow\rangle$
and $|\downarrow\rangle$, representing ``spin up'' and ``spin
down''). A general state is
\[
    |\psi\rangle = \alpha |0\rangle + \beta |1\rangle, \qquad
    |\alpha|^2 + |\beta|^2 = 1.
\]

The observable corresponding to ``spin in the $z$ direction'' is the
operator $\hat{S}_z$, which has eigenstates $|0\rangle$ and
$|1\rangle$ with eigenvalues $+\hbar/2$ and $-\hbar/2$, respectively.
If we measure $\hat{S}_z$ on the state $|\psi\rangle$:

\begin{itemize}[leftmargin=2em]
    \item The probability of obtaining $+\hbar/2$ is $|\alpha|^2$.
    \item The probability of obtaining $-\hbar/2$ is $|\beta|^2$.
    \item These probabilities sum to 1, as required.
    \item After the measurement, the state collapses to $|0\rangle$
    or $|1\rangle$ depending on the outcome.
\end{itemize}

If $\alpha = \beta = 1/\sqrt{2}$, the two outcomes are equally
probable. If $\alpha = 1, \beta = 0$, the system is already in
$|0\rangle$ and the outcome is certainly $+\hbar/2$.

\section{The measurement problem}
\label{sec:measurement}

We now have all the elements needed to state the central problem of
quantum interpretation.

\subsection{Two kinds of evolution}

The standard formalism contains two distinct kinds of dynamics.

\textbf{Schr\"odinger evolution} (Postulate 2): the wave function
evolves smoothly, deterministically, linearly, and reversibly,
according to $i\hbar \frac{d}{dt} |\psi\rangle = \hat{H} |\psi\rangle$.
This evolution preserves superposition; if the initial state was a
superposition, so is the evolved state.

\textbf{Measurement collapse} (Postulate 4): when a measurement is
performed, the wave function collapses suddenly, probabilistically,
non-linearly, and irreversibly, to one of the eigenstates of the
measured observable. This evolution destroys superposition; the
post-measurement state is a definite eigenstate.

These two kinds of evolution are mutually inconsistent. Schr\"odinger
evolution preserves superposition; measurement collapse destroys it.
A given physical interaction either preserves or destroys
superposition; it cannot do both.

\subsection{What counts as a measurement?}

The standard treatment does not specify what counts as a measurement.
Is the measurement performed by a detector? By a human observer? By
the environment? At what point in the chain --- system $\to$
detector $\to$ recording apparatus $\to$ display $\to$ retina $\to$
brain $\to$ conscious experience --- does the wave function collapse?

The standard treatment is silent. Different schools have given
different answers, none entirely satisfactory.

\begin{conceptbox}[The measurement problem]
\label{conc:measurement-problem}
The Schr\"odinger equation predicts that any physical system,
including a measurement apparatus interacting with a quantum system,
will end up in a superposition of macroscopically distinct states.
But measurements always yield definite outcomes. Either:
\begin{enumerate}[leftmargin=1.5em]
    \item The Schr\"odinger equation is incomplete, and some
    additional dynamics governs collapse; or
    \item The Schr\"odinger evolution is the only fundamental
    dynamics, and measurement collapse is some kind of effective
    description of something else; or
    \item There is no measurement collapse at all, and we are
    misdescribing what happens at measurement.
\end{enumerate}
The history of quantum interpretation is the history of attempts to
resolve this problem.
\end{conceptbox}

\section{Major interpretations of quantum mechanics}
\label{sec:interpretations}

We now survey the major interpretations of quantum mechanics. Each
takes a different position on the measurement problem and, by
extension, on the broader metaphysical questions: what is the wave
function, what is collapse, what is reality, what is the role of the
observer.

\subsection{The Copenhagen interpretation}

The Copenhagen interpretation, developed by Niels Bohr, Werner
Heisenberg, and others in the 1920s, is the textbook ``orthodox''
view from roughly 1927 to the 1980s. In its standard form:

\begin{itemize}[leftmargin=2em]
    \item The wave function is not a description of physical reality;
    it is a tool for computing measurement-outcome probabilities.
    \item The boundary between quantum system and classical
    measurement apparatus is real and primitive; one does not ask
    quantum-mechanical questions about the apparatus.
    \item Measurement collapse is a primitive process, not derivable
    from anything more fundamental; it is what happens when the
    quantum domain interacts with the classical.
    \item It is meaningless to ask what the system was doing between
    measurements.
\end{itemize}

The Copenhagen view is empirically adequate --- it reproduces all
experimental predictions of quantum mechanics --- but
philosophically uncomfortable in several ways. It leaves
``measurement'' unexplained. It draws a line between the quantum
and the classical that does not appear anywhere in the formalism. It
denies that the wave function describes reality, but is unclear what
it does describe (probabilities? Knowledge? Information?). The
Copenhagen view has been called by its critics ``shut up and
calculate'' --- a methodological injunction to use the formalism
without asking what it means.

The view's longevity owed much to its institutional dominance, the
authority of Bohr, and the practical adequacy of the formalism for
the calculations physicists wanted to do. From the 1980s onward, as
philosophical engagement with quantum foundations grew and as
alternative interpretations were developed, Copenhagen has receded
in influence.

\subsection{The Many-Worlds interpretation}

In 1957 Hugh Everett, a Princeton graduate student, proposed a
radical alternative.\footnote{Everett, H. (1957), ``Relative state
formulation of quantum mechanics,'' \emph{Reviews of Modern Physics}
29: 454--462.} \textbf{Take the Schr\"odinger equation seriously and
deny that measurement collapse occurs at all.} The wave function
never collapses. What happens at measurement is that the experimenter
and the system become entangled: the experimenter who observed
$|\text{alive}\rangle$ and the experimenter who observed
$|\text{dead}\rangle$ both exist, in different branches of a single
universal wave function. All branches are equally real.

\begin{principlebox}[The Many-Worlds Interpretation (MWI)]
The wave function of the universe evolves deterministically
according to the Schr\"odinger equation. Measurement is not a
special process; it is ordinary entanglement between a system and
an observer. After a measurement with $n$ possible outcomes, the
universal wave function contains $n$ branches, each containing a
copy of the observer who observed that outcome. All branches are
equally real.
\end{principlebox}

This sounds extravagant. The mathematical claim is, however, the
\emph{conservative} one: take the Schr\"odinger equation at face
value and don't add an extra rule for measurement. The extravagant
ontology (many worlds) is the price of mathematical conservatism.

The major technical difficulty for Many-Worlds is the so-called
\emph{Born rule problem}: if all outcomes occur in different
branches, what does it mean to say that one outcome is ``more
probable'' than another? The branches are equally real; they cannot
have different probabilities in the conventional sense. Various
attempted solutions exist (decision-theoretic derivations from
Deutsch and Wallace, branch-counting arguments, derivations from
self-locating uncertainty), but none has gained universal
acceptance.

The leading contemporary defender of Many-Worlds is Sean Carroll
(\emph{Something Deeply Hidden}, 2019). Carroll's case is that the
Schr\"odinger equation is well-tested and that any addition to it
(such as objective collapse) is unmotivated; if we take the equation
seriously, we get Many-Worlds, and the apparent extravagance is
just what taking the equation seriously costs.

\subsection{Bohmian mechanics}

Louis de Broglie proposed in 1927, and David Bohm developed in
1952, an interpretation that takes a very different approach.
Particles have definite positions at all times; the wave function is
real, but it serves as a guidance field for the particles rather
than as an exhaustive description of the system.\footnote{Bohm, D.
(1952), ``A suggested interpretation of the quantum theory in terms
of `hidden' variables, I and II,'' \emph{Physical Review} 85:
166--193.}

\begin{principlebox}[Bohmian mechanics]
\label{prin:bohmian}
The state of an $N$-particle system is given by:
\begin{enumerate}[leftmargin=1.5em]
    \item A wave function $\psi(\mathbf{q}_1, \dots, \mathbf{q}_N, t)$
    on the configuration space.
    \item The actual positions $\mathbf{Q}_1(t), \dots, \mathbf{Q}_N(t)$
    of the $N$ particles.
\end{enumerate}
The wave function evolves according to the Schr\"odinger equation.
The positions evolve according to the guidance equation:
\[
    \frac{d\mathbf{Q}_k}{dt} = \frac{\hbar}{m_k}
    \text{Im}\left( \frac{\nabla_k \psi}{\psi}
    \right)\bigg|_{\mathbf{q}_1 = \mathbf{Q}_1, \dots, \mathbf{q}_N = \mathbf{Q}_N}.
\]
Particles never lose their definite positions; the wave function never
collapses; measurement is just an interaction that produces correlations
between system and apparatus.
\end{principlebox}

Bohmian mechanics is empirically equivalent to standard quantum
mechanics --- it makes the same predictions for every experiment ---
but it has a very different ontology. Particles are real and have
definite positions. The wave function is real and guides them. The
formalism is deterministic: given the initial wave function and the
initial positions, the future positions are fully determined.

The interpretation is non-local: the guidance for one particle
depends, in general, on the positions of all the other particles,
no matter how far away. Non-locality is unavoidable in any
hidden-variable theory consistent with quantum predictions, by
Bell's theorem (which we treat in \xref{sec:bell}); Bohmian
mechanics accepts the non-locality openly.

For our project, Bohmian mechanics has a special role. It provides a
formal model in which:
\begin{itemize}[leftmargin=2em]
    \item Trajectories are fully determined.
    \item Motion is real (the particle actually moves, not just its
    description).
    \item Probability is epistemic (we don't know the initial
    position; if we did, we would predict the trajectory).
\end{itemize}
This is structurally what the classical \arb{kasb} doctrine was
trying to articulate. Chapter 16 develops the connection.

\subsection{QBism (Quantum Bayesianism)}

QBism, developed primarily by Christopher Fuchs, R\"udiger Schack,
and N. David Mermin, takes yet another approach.\footnote{Fuchs, C.
A., Mermin, N. D., \& Schack, R. (2014), ``An introduction to QBism
with an application to the locality of quantum mechanics,''
\emph{American Journal of Physics} 82: 749--754.} The wave function
does not describe an objective state of the world; it describes the
beliefs of an agent. Probability is irreducibly perspectival.

\begin{principlebox}[QBism]
\label{prin:qbism}
The wave function $|\psi\rangle$ assigned to a system by an agent
represents that agent's degrees of belief about the outcomes of
measurements they might perform. Different agents can have
different valid wave functions for the same system, depending on
what each agent knows. The Born rule is a personalist constraint
on coherent gambling; quantum probabilities are subjective
probabilities in a precise Bayesian sense.

Measurement collapse is the agent's update of beliefs in light of
new information. It is not a physical process happening in the
world; it is an epistemic event happening in the agent.
\end{principlebox}

QBism is radical in a different way from Many-Worlds. It does not
posit additional ontological structure; it removes ontological
structure that other interpretations posit. The wave function, on
QBism, does not represent anything in the world; it represents the
agent's relationship to the world. Probability is not an objective
feature of physical processes; it is the agent's own commitment.

For our project, QBism has special importance because it gives us
the conceptual structure in which probability can be irreducibly
perspectival. Different agents (or knowledge states) can have
different probability assignments without contradiction. Applied to
the question of \arb{qadar}: divine knowledge is perfect (the divine
``wave function'' is, in QBism's language, fully concentrated on
the actual outcome); human knowledge is finite (the human wave
function spreads probability over multiple outcomes). Both are
correct; the difference is one of perspective. There is no
contradiction. We develop this in Chapter 16.

\subsection{Objective collapse theories}

Objective collapse theories, developed in various forms by
Ghirardi, Rimini, and Weber (GRW), Pearle, Diosi, and Penrose,
modify the Schr\"odinger equation by adding a small additional term
that produces spontaneous collapse at random times. The collapse is
a real physical process, not a primitive of measurement.

The most physically interesting version is Penrose's gravity-induced
collapse, in which the collapse rate is determined by the
gravitational self-energy of the superposition. Macroscopic objects
collapse rapidly (because their gravitational self-energies are
large); microscopic objects collapse slowly. The transition between
quantum and classical regimes is then explained by the same
underlying mechanism, with no need for a primitive
quantum/classical boundary.

Objective collapse theories are empirically distinguishable from
standard quantum mechanics in principle, but the predicted effects
are extremely small and have not yet been definitively detected or
ruled out. The Penrose-Hameroff Orch-OR theory of consciousness
(treated in Chapter 8) is built on objective collapse.

\subsection{Relational quantum mechanics}

Carlo Rovelli's relational quantum mechanics (RQM) takes the
position that the state of a system is always defined relative to
another system. There is no absolute, observer-independent state.
Different observers can describe the same system differently, and
both descriptions are correct, because state is inherently
relational.

RQM is in some ways a synthesis of QBism's perspectival approach
and Bohmian-style realism: states are real, but they are real
relative to other systems, not absolutely. The position has gained
attention in the philosophical literature but is less developed
formally than the major alternatives.

\subsection{Comparing the interpretations}

The interpretations differ on several dimensions:

\begin{center}
\begin{tabular}{l@{\hspace{1em}}c@{\hspace{1em}}c@{\hspace{1em}}c@{\hspace{1em}}c}
\toprule
\textbf{Interpretation} & \textbf{Wave fn.\ real?} & \textbf{Determ.?} & \textbf{Local?} & \textbf{Collapse?} \\
\midrule
Copenhagen & no & no & yes (?) & primitive \\
Many-Worlds & yes & yes & yes & no \\
Bohmian & yes & yes & no & no \\
QBism & no & no (?) & yes & epistemic \\
Objective collapse & yes & no & yes & physical \\
Relational QM & relative & --- & --- & relative \\
\bottomrule
\end{tabular}
\end{center}

The choice among interpretations is not adjudicable by current
experiments; they all reproduce the same empirical predictions. The
choice is at the level of philosophical commitment, informed but
not determined by the physics.

\section{Bell's theorem and the question of locality}
\label{sec:bell}

A brief treatment of one of the most important results in quantum
foundations.

\subsection{The EPR argument}

In 1935, Einstein, Podolsky, and Rosen (EPR) argued that quantum
mechanics is incomplete.\footnote{Einstein, A., Podolsky, B., \&
Rosen, N. (1935), ``Can quantum-mechanical description of physical
reality be considered complete?'' \emph{Physical Review} 47:
777--780.} Their argument used a thought experiment with two
particles whose properties are correlated; the wave function for the
pair is a superposition such that measuring one particle yields
information about the other.

If, said EPR, the second particle has a definite property (which we
can determine without disturbing it, by measuring the first), then
that property must have been definite all along. But the wave
function does not assign it a definite value. Therefore the wave
function does not contain all the relevant information; quantum
mechanics is incomplete. There must be ``hidden variables'' that
the formalism omits.

\subsection{Bell's theorem}

In 1964, John Bell proved a remarkable theorem.\footnote{Bell, J. S.
(1964), ``On the Einstein-Podolsky-Rosen paradox,'' \emph{Physics}
1: 195--200.} Any local hidden-variable theory --- any theory in
which physical properties are determined by hidden variables and in
which no physical influence propagates faster than light --- makes
predictions that differ from those of quantum mechanics. Specifically,
Bell derived an inequality that any local hidden-variable theory must
satisfy and that quantum mechanics violates.

Subsequent experiments (Aspect 1982, Hensen et al.\ 2015, and many
others) have confirmed the quantum predictions to high precision and
violated Bell's inequality. The conclusion is that no local
hidden-variable theory is consistent with the experimental facts.

\begin{conceptbox}[What Bell's theorem rules out]
\label{conc:bell}
Any theory satisfying \emph{both} of the following conditions is
ruled out by experiment:
\begin{enumerate}[leftmargin=1.5em]
    \item \emph{Locality}: no influence propagates faster than
    light, and physical properties at one location depend only on
    causal influences from the past light cone.
    \item \emph{Hidden variables}: physical properties are determined
    by underlying variables, possibly inaccessible to measurement,
    that fully specify the state.
\end{enumerate}
At least one of these must be wrong.
\end{conceptbox}

The implications are striking. If we want a hidden-variable theory
(such as Bohmian mechanics), we must accept non-locality. If we
want locality, we must give up the idea that quantum systems have
definite properties prior to measurement. Either way, the world is
not the local-realist Newtonian universe that classical physics
described.

\subsection{Why this matters for our project}

Bell's theorem is important to the present project for two reasons.

First, it shows that the classical Newtonian-realist worldview is
empirically refuted. The world is not a collection of objects with
definite properties interacting locally according to deterministic
laws. Some other framework is required. This dissolves the
philosophical pressure that made Islamic metaphysical claims look
implausible from a classical-materialist perspective.

Second, it makes the Bohmian and QBist interpretations
philosophically defensible. Bohmian mechanics requires non-locality
but is consistent with Bell; QBism gives up the assumption that
quantum systems have definite properties prior to measurement and
is also consistent with Bell. Both interpretations, which we will
use in Chapter 16 to articulate \arb{qadar}, are not refuted by
Bell; in fact, they are the natural responses to it.

\section{Implications for the Hidden Architecture project}
\label{sec:qm-implications}

We have developed enough quantum mechanics to draw the implications
that the rest of the book will depend on. Three implications are
central.

\subsection{Determinism and motion are not opposed}

The classical \arb{kalam} debate on \arb{qadar} assumed that
determinism (everything is determined by divine decree) and motion
(human action is real) were in tension. The Bohmian interpretation
shows that this assumption was wrong. In Bohmian mechanics, the
trajectory of a particle is fully determined by the initial wave
function and initial position, yet the motion is real and the
particle actually traverses the trajectory.

Translated to the \arb{qadar} question: the trajectory of a human
life is fully determined by divine decree (analogous to the wave
function), yet the action is real and the person actually performs
the action (analogous to the particle's motion along the trajectory).
There is no contradiction.

This is not to claim that the human mind is governed by Bohmian
quantum mechanics. The claim is structural: Bohmian mechanics shows
that the structure ``determined trajectory + real motion'' is
mathematically coherent. Whatever the actual mechanism of human
action is, it can have this structure without contradiction. The
classical theological worry that determinism makes motion illusory
was a worry about a structure that turns out to be available.

We develop this in detail in Chapter 16.

\subsection{Probability can be perspectival}

The classical theological worry that divine perfect knowledge
contradicts genuine human uncertainty assumed that probability must
be a single objective fact. QBism shows that this assumption was
wrong. Probability can be irreducibly perspectival: different
agents, with different information, can have different valid
probability assignments for the same event without contradiction.

Translated: divine knowledge of all events is perfect. The divine
``probability assignment'' is concentrated on the actual outcome.
Human knowledge is finite. The human ``probability assignment'' is
spread over multiple outcomes. Both are correct; they reflect the
different epistemic positions of the divine and human perspectives.

This too is structural. The claim is not that the divine mind is a
QBist agent. The claim is that the structure ``different valid
probability assignments for the same event'' is mathematically and
philosophically coherent, as QBism shows for the quantum case.

\subsection{The materialist worldview is not forced}

Classical materialism --- the view that reality consists of
particles in definite positions interacting locally according to
deterministic laws --- is not the worldview of contemporary
fundamental physics. Bell's theorem rules out local realism. Some
form of either non-locality (Bohmian) or non-realism (QBism) is
required.

This does not establish Islamic metaphysics. It does, however,
remove an obstacle. The reader who comes to this book convinced
that ``science has shown that materialism is correct and that
metaphysical claims about decree, intention, and blessing are
therefore unscientific'' is operating with a false picture of where
science currently stands. The materialist picture is one option;
several others are equally consistent with the empirical evidence;
each makes different metaphysical commitments. The choice among
them is philosophical, not empirical.

The Islamic metaphysical claims are not in conflict with science.
They are in conflict with one philosophical interpretation of
science --- an interpretation that has been declining in
professional credibility for several decades.

\begin{summarybox}[Chapter summary]
\textbf{Mathematical preliminaries}: complex numbers (with modulus
and conjugate), complex vector spaces, inner products, Hilbert
spaces, Dirac bra-ket notation.

\textbf{The four postulates of quantum mechanics}:
\begin{itemize}[leftmargin=1.5em]
    \item States are normalized vectors in a complex Hilbert space.
    \item Time evolution is given by the Schr\"odinger equation;
    deterministic, linear, reversible.
    \item Observables are Hermitian operators; their eigenvalues
    are real and possible measurement outcomes.
    \item The Born rule: probability of outcome $a$ is
    $|\langle a | \psi\rangle|^2$; after measurement, state
    collapses to $|a\rangle$.
\end{itemize}

\textbf{The measurement problem}: Schr\"odinger evolution preserves
superposition; measurement collapse destroys it. The two cannot
both be fundamental. Different interpretations resolve this
differently.

\textbf{Major interpretations}:
\begin{itemize}[leftmargin=1.5em]
    \item Copenhagen: collapse is primitive; wave function is a
    calculation tool.
    \item Many-Worlds: no collapse; all outcomes occur in
    different branches.
    \item Bohmian: particles have definite positions; wave function
    guides them; deterministic; non-local.
    \item QBism: wave function represents agent beliefs; probability
    is perspectival; collapse is epistemic.
    \item Objective collapse: collapse is a real physical process,
    via additional dynamics.
\end{itemize}

\textbf{Bell's theorem}: rules out local hidden-variable theories.
The world is either non-local (Bohmian) or has no definite
properties prior to measurement (QBism, Many-Worlds).

\textbf{For the Hidden Architecture project}: Bohmian mechanics
provides the formal structure for ``determined trajectory + real
motion'' that the \arb{kasb} doctrine was attempting; QBism
provides the structure for ``perspectival probability'' that
resolves the apparent tension between divine knowledge and human
uncertainty; classical materialism is not the only option, removing
the philosophical obstacle to the metaphysical claims.

\textbf{What comes next}: Chapter 6 develops ergodicity economics,
the second major resource of Part II.
\end{summarybox}

\cleardoublepage
